% Figure: the air-standard Brayton cycle on temperature-entropy axes.
% Two constant-pressure legs (heat addition, heat rejection) joined by two
% isentropic legs (compression 1->2, expansion 3->4). Drawn schematically;
% the constant-pressure curves diverge with increasing entropy.
\begin{figure}[htbp]
\centering
\begin{tikzpicture}[font=\small, scale=1.0]
  \begin{axis}[
      width=10cm, height=7cm,
      xlabel={specific entropy $s$ (kJ/(kg$\cdot$K))},
      ylabel={temperature $T$ (K)},
      xmin=6.4, xmax=8.2, ymin=250, ymax=1600,
      axis lines=left,
      xtick=\empty,
      ytick={300,600,900,1200,1500},
      tick label style={font=\scriptsize},
      label style={font=\scriptsize},
      clip=false,
    ]
    % low-pressure isobar (rejection line), p1: T = 290 * exp((s - 6.60)/cp)
    \addplot[engblue, thick, domain=6.60:8.05, samples=60] {290*exp((x-6.60)/1.005)};
    % high-pressure isobar (addition line), p2: shifted left in s by R ln(rp)
    \addplot[engred, thick, domain=6.60:8.05, samples=60] {576*exp((x-7.06)/1.005)};
    % isentropic compression 1 -> 2 (vertical, constant s)
    \draw[engblack, thick] (axis cs:6.60,290) -- (axis cs:6.60,576);
    % isentropic expansion 3 -> 4 (vertical, constant s)
    \draw[engblack, thick] (axis cs:7.50,1450) -- (axis cs:7.50,730);
    % state markers
    \node[circle, fill=engblack, inner sep=1.3pt, label={below left:1}] at (axis cs:6.60,290) {};
    \node[circle, fill=engblack, inner sep=1.3pt, label={above left:2}] at (axis cs:6.60,576) {};
    \node[circle, fill=engblack, inner sep=1.3pt, label={above right:3}] at (axis cs:7.50,1450) {};
    \node[circle, fill=engblack, inner sep=1.3pt, label={below right:4}] at (axis cs:7.50,730) {};
    % leg labels
    \node[engred, font=\scriptsize, anchor=south] at (axis cs:7.05,1180) {heat in (combustor)};
    \node[engblue, font=\scriptsize, anchor=north] at (axis cs:7.05,420) {heat out (exhaust)};
    \node[enggray, font=\scriptsize, rotate=90, anchor=south] at (axis cs:6.55,440) {compressor};
    \node[enggray, font=\scriptsize, rotate=90, anchor=north] at (axis cs:7.55,1090) {turbine};
  \end{axis}
\end{tikzpicture}
\caption[The air-standard Brayton cycle on temperature-entropy axes]{The
air-standard Brayton cycle, the model for a gas turbine, drawn on
temperature-entropy axes. Air is compressed isentropically from state 1 to
state 2, heated at constant pressure in the combustor from 2 to 3, expanded
isentropically through the turbine from 3 to 4, and rejected at constant
pressure from 4 back to 1 (in an open engine the rejection happens in the
atmosphere rather than in a closed loop). The two isobars diverge as entropy
rises, so the heat added at the high temperatures of the combustor exceeds the
heat rejected at the low temperatures of the exhaust, and the enclosed area is
the net work. Raising the turbine inlet temperature at state 3 stretches the
cycle upward and is the lever that has driven gas-turbine efficiency for
decades, bounded by the metallurgy and the cooling of the first-stage blades.}
\label{fig:vol04ch01:brayton-ts}
\end{figure}
